\documentclass[11pt]{amsart}
\usepackage{amsmath,amssymb,amsthm,mathtools}
\usepackage[margin=1in]{geometry}
\usepackage{hyperref}
\usepackage{listings}
\lstset{basicstyle=\ttfamily\footnotesize,breaklines=true,frame=single,
  showstringspaces=false,columns=fullflexible,keepspaces=true}

\newtheorem{theorem}{Theorem}[section]
\newtheorem{proposition}[theorem]{Proposition}
\newtheorem{lemma}[theorem]{Lemma}
\newtheorem{corollary}[theorem]{Corollary}
\theoremstyle{definition}
\newtheorem{definition}[theorem]{Definition}
\theoremstyle{remark}
\newtheorem{remark}[theorem]{Remark}

\newcommand{\RH}{\mathrm{RH}}
\newcommand{\HW}{\mathcal H_W}
\newcommand{\T}{\mathcal T}
\newcommand{\Tc}{\mathcal T_{\mathrm{cont}}}
\newcommand{\Ainf}{A_\infty}
\newcommand{\sltwo}{\mathfrak{sl}_2}
\DeclareMathOperator{\spec}{spec}
\DeclareMathOperator{\sgn}{sgn}
\DeclareMathOperator{\Real}{Re}

\title[The Lefschetz Dichotomy]{The Lefschetz dichotomy: why the hard-Lefschetz route to the
Riemann Hypothesis is obstructed by linear independence}
\author{David Alejandro Trejo Pizzo}
\thanks{Independent Researcher. \texttt{dtrejopizzo@gmail.com}}


\begin{document}
\nocite{bgWeil49,bgDeligne,bgConnes99,bgCC10,bgDeninger,bgBK,bgMont73,bgKS,bgBombieri00,bgTitchmarsh,bgEdwards,bgIK}
\begin{abstract}
Weil's proof of the Riemann Hypothesis for curves over finite fields is an
intersection form plus a Hodge index: positivity on the primitive cohomology,
delivered by an ample class (a hard-Lefschetz $\sltwo$-action). Transcribing the
construction to $\operatorname{Spec}\mathbb Z$ produces, unconditionally, every
ingredient except that ample class. We make the missing object precise and prove
that it cannot be supplied along this route. First we record an
\emph{independence filter}: any raising operator built as a function of the
Frobenius flow $\T$ (whose spectrum is the zeros) only restates the Weil
positivity. Then we construct the candidate Kähler data from the archimedean
place: the density $\Psi(r)=\Real\psi(\tfrac14+\tfrac{ir}{2})-\log\pi$ is the
smooth zero-counting density; it is strictly increasing on $r>0$ with a unique
sign change $r_0=6.2898\ldots$ lying below the first zero, hence positive on the
spectrum but indefinite on the low band. We show the resulting form is indefinite
on the primitive part, exhibit an explicit hard-Lefschetz $\sltwo$ on the
archimedean continuum whose raising operator is provably outside the von Neumann
algebra of $\T$ --- the first such operator this line of work produces --- and verify
that the transport to the discrete zeros is exactly the explicit formula, with
the Weil positivity emerging only as the residual of an almost total
archimedean--prime cancellation. Finally we prove the \emph{Lefschetz dichotomy}:
on the zeros the linear independence of the ordinates forbids any Tate-twist
raising operator, while on the continuum the only available $\sltwo$ is principal
series and carries no Hodge--Riemann form. Thus an arithmetic spectrum and a
Hodge--Riemann-carrying grading are linearly-independence-incompatible. This is a
theorem about why $\zeta$ resists the Lefschetz route, sharpening the
finitization-under-LI obstruction of the series. All claims are accompanied by
reproducible computations; none asserts a proof of RH.
\end{abstract}
\maketitle

\tableofcontents
\newpage

\section{Introduction}
Weil's proof of $\RH$ for a curve $C/\mathbb F_q$ has two parts: an intersection
pairing on the surface $C\times C$, and the Hodge index theorem, which makes the
pairing negative definite on the primitive part. The second part is supplied by
an \emph{ample class} on the surface --- equivalently a hard-Lefschetz
$\sltwo$-action $(L,\Lambda,H)$ whose positivity is the Hodge--Riemann bilinear
relation. The Frobenius acts on $H^1$, whose eigenvalues are the zeros, and the
ample class is independent geometric input.

The program of which this paper is a part has built, unconditionally and with
numerical verification, the analogue of the \emph{first} part for
$\operatorname{Spec}\mathbb Z$: a Kre\u\i n space $\HW$ with an intersection
pairing whose negative index counts the off-line zeros, a Lefschetz trace
identity (the explicit formula), a Rankin--Selberg effectivity, and a positive
ample cone~\cite{RP8,RP16,RP17}. The Riemann Hypothesis is then equivalent to the
\emph{second} part: positivity of the pairing on the primitive subspace --- the
arithmetic Hodge index. Eight successive sharpenings reduced the gap to a single
missing object: an independent, positive, hard-Lefschetz $\sltwo$-action (a Kähler
class $\omega$) on the primitive cohomology, commuting appropriately with the
Frobenius flow $\T$ whose spectrum is the set of ordinates $\{\gamma_\rho\}$.

This paper does two things. It records why the obvious construction --- building
$\omega$ from the spectral data $(\T,J,Q)$ --- can only restate $\RH$ (\S2). And
it carries out the one construction that escapes this circularity, the archimedean
one, to its conclusion (\S\S3--6), proving along the way the first
hard-Lefschetz raising operator independent of $\T$ (\S5), and ending at a clean
structural obstruction: the \emph{Lefschetz dichotomy} (\S7). The dichotomy
explains, in representation-theoretic terms, why the linear independence of the
zeros (the hypothesis $\mathrm{LI}$, that the positive ordinates are
$\mathbb Q$-linearly independent) blocks the Lefschetz route. Nothing here proves
or assumes $\RH$; the result is negative and structural, and every quantitative
claim is reproducible from the code in Appendix~\ref{app:code}.

\section{The independence filter}
We work on the primitive part $\HW^{\mathrm{prim}}$, an infinite-dimensional real
Kre\u\i n space carrying: the indefinite Weil pairing $Q$ (negative index $=$
number of off-line zeros), a self-adjoint Frobenius operator $\T$ with
$\spec(\T)=\{\gamma_\rho\}$, and a complex structure $J$ from the spectral
splitting~\cite{RP8}. Let $W^*(\T)$ denote the von Neumann algebra generated by
$\T$.

\begin{proposition}[Determinacy / the filter]\label{prop:filter}
The only $(1,1)$-form compatible with $(J,Q)$ in the standard way is
$\omega(x,y)=Q(Jx,y)$, and its positivity $\omega(x,Jx)=Q(Jx,Jx)>0$ is exactly
$\RH$. Consequently, if a raising operator $L$ lies in $W^*(\T)$ --- i.e.\
$L=f(\T)$ for some Borel $f$ --- then the associated form $Q(Jf(\T)\cdot,\cdot)$ is
a function of the spectral data and asserting its positivity restates $\RH$. An
$\omega$ that \emph{forces} positivity must use a raising operator
$L\notin W^*(\T)$.
\end{proposition}

\begin{proof}
Compatibility of a $(1,1)$-form with $(J,Q)$ means $\omega(x,y)=Q(Jx,y)$; this
fixes $\omega$. Its positivity is the second Riemann bilinear relation
$Q(Jx,Jx)>0$, equivalent by the Weil criterion to the vanishing of the negative
index, i.e.\ $\RH$. If $L=f(\T)$ then $L$ commutes with $\T$ and is diagonal in
the spectral resolution of $\T$; every quantity built from $L,J,Q$ is a function
of that resolution, so its sign is the sign of $Q$ on the spectral subspaces ---
the statement $\RH$ itself.
\end{proof}

The filter is the lesson of the eight attempts in compact form: a candidate Kähler
class is genuine input only if its raising operator is \emph{not} a function of
the Frobenius. We now produce one.

\section{The archimedean density}
The archimedean place contributes to the Weil quadratic form, for a test function
$f$ with Fourier transform $\widehat f$, the term
\begin{equation}\label{eq:Ainf}
\Ainf(f,f)=\frac1{2\pi}\int_{-\infty}^{\infty}|\widehat f(r)|^2\,\Psi(r)\,dr,
\qquad
\Psi(r)=\Real\psi\!\left(\tfrac14+\tfrac{ir}{2}\right)-\log\pi,
\end{equation}
where $\psi=\Gamma'/\Gamma$ and $\Psi$ is the smooth density of the archimedean
factor $\Gamma_{\mathbb R}(s)=\pi^{-s/2}\Gamma(s/2)$ in the explicit formula. The
following structure is elementary but decisive.

\begin{proposition}[Structure of $\Psi$]\label{prop:psi}
$\Psi$ is even and real-analytic, and:
\begin{enumerate}
\item[\rm(i)] $\Psi'(r)=\displaystyle\sum_{n\ge0}\frac{(\tfrac14+n)\,r}
{\bigl((\tfrac14+n)^2+(r/2)^2\bigr)^2}>0$ for $r>0$ (strictly increasing);
\item[\rm(ii)] $\Psi$ has a unique positive zero $r_0=6.2898359888\ldots$, global
minimum $\Psi(0)=\psi(\tfrac14)-\log\pi=-\gamma-3\log2-\tfrac{\pi}{2}-\log\pi
=-5.37218\ldots$;
\item[\rm(iii)] $\Psi(r)=\log\dfrac{|r|}{2\pi}+o(1)$, and
$\tfrac1{2\pi}\Psi(r)$ equals the Riemann--von Mangoldt zero density $N'(r)$ to
leading order.
\end{enumerate}
In particular $\Psi(r)>0\iff|r|>r_0$; the only negativity is the bounded band
$|r|<r_0$, which lies strictly inside the zero-free gap, $r_0<\gamma_1
=14.1347\ldots$, so $\Psi>0$ on the entire support of the zeros.
\end{proposition}

\begin{proof}
Evenness and analyticity are clear. For (i), differentiate \eqref{eq:Ainf}:
$\Psi'(r)=\Real\bigl[\tfrac{i}{2}\psi'(\tfrac14+\tfrac{ir}{2})\bigr]
=-\tfrac12\Im\psi'(\tfrac14+\tfrac{ir}{2})$, and
$\psi'(z)=\sum_{n\ge0}(z+n)^{-2}$ gives
$\Im\psi'(\tfrac14+\tfrac{ir}{2})=-\sum_n\frac{2(\tfrac14+n)(r/2)}
{((\tfrac14+n)^2+(r/2)^2)^2}<0$ for $r>0$; the stated series follows. Hence
$\Psi$ increases from $\Psi(0)<0$ to $+\infty$, crossing zero once at $r_0$;
(ii)--(iii) are the standard values and the digamma asymptotics
$\Real\psi(\tfrac14+\tfrac{ir}{2})=\log\tfrac{|r|}{2}+o(1)$. The numerical
constants and $r_0<\gamma_1$ are verified to $40$ digits in
Appendix~\ref{app:code}.
\end{proof}

So the archimedean form is a strictly positive metric on the spectral bulk
($|r|\ge\gamma_1$, where $\T$ has spectrum) and indefinite only on a finite band
in the spectral gap.

\section{Indefiniteness on the primitive part}
A natural hope is that the band $|r|<r_0$ is absorbed by the ample/pole part of
the construction (the de Vallée Poussin edge positivity of~\cite{RP16}). It is
not.

\begin{proposition}[$\Ainf$ is indefinite on primitives]\label{prop:indef}
The vectors $\widehat f_\sigma(r)=(r^2+\tfrac14)e^{-r^2/2\sigma^2}$ are primitive
($\widehat f_\sigma(\pm\tfrac{i}{2})=0$, so the pole/ample term vanishes on them),
and $\Ainf(f_\sigma,f_\sigma)<0$ for small $\sigma$: numerically
$\Ainf/\|f_\sigma\|^2=-1.82$ at $\sigma=1$, crossing to positive near
$\sigma\approx5$. Hence $\Ainf$ is indefinite on $\HW^{\mathrm{prim}}$ and the
band is not absorbed by the pole part.
\end{proposition}

\begin{proof}
$(i/2)^2+\tfrac14=0$ gives $\widehat f_\sigma(\pm i/2)=0$, so $f_\sigma$ is
primitive and the pole term $2|\widehat f_\sigma(i/2)|^2$ vanishes identically.
For small $\sigma$ the mass of $|\widehat f_\sigma|^2$ concentrates in
$|r|<r_0$ where $\Psi<0$, giving $\Ainf<0$; the values are computed in
Appendix~\ref{app:code}.
\end{proof}

The signed form factors through a positive metric and a grading sign:
$\Ainf=g_\infty(S\,\cdot,\cdot)$ with $g_\infty=|\Psi|(\Tc)\succ0$ and
$S=\sgn(\Tc^2-r_0^2)\in\{\pm1\}$ --- the shape of a Hodge--Riemann form, a positive
metric times a degree sign. The indefiniteness is a candidate grading, not a
defect; whether it organizes into a hard-Lefschetz $\sltwo$ is the next question.

\section{The continuum $\sltwo$ and the discrete no-go}
The Frobenius $\T$ acts on $\HW$ with simple spectrum, so its commutant is a
maximal abelian subalgebra (a MASA): any $L$ commuting with $\T$ is a function of
$\T$ and, by Proposition~\ref{prop:filter}, useless. The correct hard-Lefschetz
condition is therefore not commutation but the \emph{Tate twist}
$[\T,L]=cL$ --- the analogue of $\mathrm{Frob}\cdot L=q\,L\cdot\mathrm{Frob}$ on
$C\times C$, where $L$ raises the weight and scales the Frobenius eigenvalue.

\begin{proposition}[No Tate-twist Lefschetz on the zeros]\label{prop:nogo}
If $[\T,L]=cL$ with $c\ne0$ and $L\ne0$, then $L$ maps the $\gamma$-eigenspace of
$\T$ into the $(\gamma+c)$-eigenspace, so the zero set must satisfy
$\{\gamma_\rho\}+c\subseteq\{\gamma_\rho\}$, i.e.\ be invariant under translation
by $c$. No $c\ne0$ has this property: a single inclusion $\gamma,\gamma+c,
\gamma+2c\in\{\gamma_\rho\}$ already is a nontrivial $\mathbb Q$-linear relation
among ordinates, excluded by $\mathrm{LI}$ \cite{RP15}; unconditionally, the
variable spacing of the zeros rules out exact translation invariance. Hence no
hard-Lefschetz raising operator exists on $\HW$ itself.
\end{proposition}

The remedy is to pass to the archimedean continuum $L^2(\mathbb R,dr)$, where the
multiplication operator $\Tc$ has spectrum all of $\mathbb R$ --- which \emph{is}
translation invariant.

\begin{proposition}[Explicit continuum $\sltwo$]\label{prop:sl2}
On $L^2(\mathbb R,dr)$ put $L=T_c$ (translation by $c$),
$H=\tfrac2c\,\Tc$, and $\Lambda=T_{-c}\,m(\Tc)$ with $m(r)=-\tfrac{r^2}{c^2}
+\tfrac{r}{c}$. Then
\[
[\Tc,L]=cL,\qquad [H,L]=2L,\qquad [H,\Lambda]=-2\Lambda,\qquad [L,\Lambda]=H,
\]
so $(L,\Lambda,H)$ is an $\sltwo$-triple, Tate-twist compatible with the Frobenius
flow. The raising operator $L=T_c$ is not a multiplication operator, hence
$L\notin W^*(\Tc)$: the independence filter of Proposition~\ref{prop:filter} is
cleared by an explicit operator. The grading $H=\tfrac2c\Tc$ has continuous
spectrum --- a ``Frobenius-as-a-flow'' Lefschetz, not an integer grading.
\end{proposition}

\begin{proof}
$[\Tc,T_c]=cT_c$ is the canonical commutation of multiplication and translation.
Then $[H,L]=\tfrac2c[\Tc,T_c]=2T_c=2L$. With
$(T_{-c}\,m\,T_c f)(r)=m(r+c)f(r)$ one computes
$[L,\Lambda]=[T_c,T_{-c}m]=m(r)-m(r+c)=\tfrac2c r=H$ (using
$m(r)-m(r+c)=-\tfrac{1}{c^2}(r^2-(r+c)^2)+\tfrac1c(r-(r+c))=\tfrac2c r$), and
$[H,\Lambda]=\tfrac2c[\Tc,T_{-c}]m=-2\Lambda$. The relations are verified
numerically to $10^{-12}$ in Appendix~\ref{app:code}. A translation is not
diagonal in the position representation, so $L\notin W^*(\Tc)$.
\end{proof}

Propositions~\ref{prop:nogo}--\ref{prop:sl2} are the two horns to be joined in
\S7. First we check what the construction buys downstairs.

\section{The transport is the explicit formula}
The continuum structure connects to the arithmetic zeros only through sampling.

\begin{proposition}[Transport identity]\label{prop:transport}
For primitive $f$ the explicit formula reads
\begin{equation}\label{eq:transport}
Q(f,f)=\sum_\rho |\widehat f(\gamma_\rho)|^2
=\Ainf(f,f)-P(f,f),\qquad
P(f,f)=2\sum_{n\ge2}\frac{\Lambda(n)}{\sqrt n}\,g(\log n),
\end{equation}
where $g$ is the Fourier transform of $|\widehat f|^2$ (the pole term drops since
$\widehat f(\pm i/2)=0$). On the test family of Proposition~\ref{prop:indef} the
identity holds to $10^{-12}$--$10^{-20}$, and $\Ainf$ and $P$ are large, negative,
and nearly equal, so $Q$ is the small non-negative residual of an almost total
archimedean--prime cancellation.
\end{proposition}

Table~\ref{tab:transport} records the computation (first $300$ zeros, von
Mangoldt to $2\cdot10^4$). The point is structural: the discrete positivity
$Q\ge0$ is \emph{not} the transported image of the upstairs metric --- on these
primitives $\Ainf<0$ --- but the residue $\Ainf-P$, the classical wrong-sign
obstruction~\cite{RP13}.

\begin{table}[h]
\centering
\begin{tabular}{c|rrrr}
$\sigma$ & $Q=\sum_\rho|\widehat f(\gamma_\rho)|^2$ & $\Ainf$ & $P$ & $Q-(\Ainf-P)$\\\hline
$2$ & $1.6\times10^{-17}$ & $-6.2225$ & $-6.2225$ & $-1.4\times10^{-20}$\\
$3$ & $1.83\times10^{-5}$ & $-21.677$ & $-21.677$ & $-2\times10^{-12}$\\
$4$ & $0.3022$ & $-24.239$ & $-24.542$ & $-3\times10^{-4}$
\end{tabular}
\caption{The transport identity \eqref{eq:transport} and the near-cancellation.}
\label{tab:transport}
\end{table}

\section{The Lefschetz dichotomy}
The continuum $\sltwo$ of Proposition~\ref{prop:sl2} clears the independence
filter, but it does not carry the positivity that hard Lefschetz needs.

\begin{proposition}[No Hodge--Riemann form on the continuum $\sltwo$]\label{prop:ps}
The representation of Proposition~\ref{prop:sl2} is principal series:
$\spec(H)=\tfrac2c\mathbb R$ is unbounded below, and the lowest-weight equation
$\Lambda v=0$ forces $m(r)v(r)=0$, i.e.\ $v$ supported on the measure-zero set
$\{0,c\}$, so there is no normalizable lowest-weight vector and no primitive
decomposition $V=\bigoplus_k L^k(\ker\Lambda)$. Hodge--Riemann bilinear relations
are defined only on lowest-weight (equivalently finite-dimensional or
discrete-series, integer-graded) $\sltwo$-representations; the continuum $\sltwo$
carries none. Concretely, the only manifestly positive form the data yield is
$g_\infty(f,f)=\tfrac1{2\pi}\int|\widehat f|^2|\Psi|$, which differs from the Weil
form $Q$ by factors $10^2$--$10^{18}$ on the test family
(Appendix~\ref{app:code}); the degree grading produces $g_\infty$, not $Q$.
\end{proposition}

\begin{theorem}[The Lefschetz dichotomy]\label{thm:dichotomy}
Under $\mathrm{LI}$ there is no hard-Lefschetz $\sltwo$-action carrying a
Hodge--Riemann form that acts on the arithmetic zero-spectrum:
\[
\boxed{\;
\begin{array}{ll}
\textbf{discrete horn:} & \text{on }\{\gamma_\rho\}\text{ no Tate-twist raising
operator exists (Prop.~\ref{prop:nogo}),}\\
\textbf{continuum horn:} & \text{on }L^2(\mathbb R,dr)\text{ the raising operator
exists but the }\sltwo\\
& \text{is principal series, with no Hodge--Riemann form
(Prop.~\ref{prop:ps}).}
\end{array}\;}
\]
Equivalently: a Hodge--Riemann-carrying grading requires lowest weight (integer
spacing), which forces the Frobenius eigenvalues into a single weight; $\mathrm{LI}$
is exactly the statement that the zeros admit no such integer-weight structure.
An arithmetic spectrum and a Hodge--Riemann grading are
$\mathrm{LI}$-incompatible.
\end{theorem}

\begin{proof}
The two horns are Propositions~\ref{prop:nogo} and~\ref{prop:ps}. For the
equivalence: in the function-field case $H^1$ is pure of weight $1$, the
Frobenius eigenvalues are Weil numbers of that single weight, and the
Lefschetz $\sltwo$ on $H^0\oplus H^1\oplus H^2$ is lowest weight, so
Hodge--Riemann applies and yields $\RH$. A lowest-weight (integer-graded)
$\sltwo$ commuting-up-to-Tate-twist with $\T$ would, by the argument of
Proposition~\ref{prop:nogo}, place the eigenvalues of $\T$ in an arithmetic
progression of weights --- a $\mathbb Q$-linear structure on $\{\gamma_\rho\}$
excluded by $\mathrm{LI}$. Hence no such grading exists on the zeros; buying one
on the continuum costs the lowest-weight property (Proposition~\ref{prop:ps}).
\end{proof}

\begin{remark}
Theorem~\ref{thm:dichotomy} refines the finitization-under-LI obstruction
of~\cite{RP15} from ``no finite model'' to a representation-theoretic statement
about gradings, and it is reached constructively: every ingredient of Weil's
scaffold is built (pairing, trace identity, effectivity, ample cone, the
zero-carrying cohomology, the Frobenius flow, and an explicit $\T$-independent
$\sltwo$), and the obstruction is isolated as the incompatibility of two
specific structures. It does not assert that $\RH$ is false, nor that no proof
exists --- only that the hard-Lefschetz mechanism, as it operates in the
function-field case, cannot be transported to $\zeta$ along this route.
\end{remark}

\section{Conclusion}
The archimedean--modular construction produces three genuine gains: a positive
metric $g_\infty=|\Psi|$ on the spectral bulk (Proposition~\ref{prop:psi}); the
first explicit hard-Lefschetz raising operator independent of the Frobenius
(Proposition~\ref{prop:sl2}); and the exact transport identity
(Proposition~\ref{prop:transport}). It does not cross the wall: the Weil
positivity is the near-cancellation residual $\Ainf-P$, and the
Lefschetz dichotomy (Theorem~\ref{thm:dichotomy}) names precisely why the
positivity cannot be supplied by a hard-Lefschetz class --- linear independence
of the zeros forbids the lowest-weight grading that Hodge--Riemann requires. The
result is a structural theorem about the obstruction, consistent with the
program's wrong-sign map~\cite{RP13,RP15}, reached this time from the
Hodge-theoretic side.

\appendix
\section{Reproducible computations}\label{app:code}
All values are produced by the following self-contained scripts (Python with
\texttt{mpmath}/\texttt{numpy}); outputs are quoted in the text.

\subsection*{A.1 The archimedean density (Propositions~\ref{prop:psi})}
\begin{lstlisting}
import mpmath as mp
mp.mp.dps = 40
logpi = mp.log(mp.pi)
Psi  = lambda r: mp.re(mp.digamma(mp.mpf('0.25')+1j*mp.mpf(r)/2)) - logpi
dPsi = lambda r: mp.re(0.5*1j*mp.polygamma(1, mp.mpf('0.25')+1j*mp.mpf(r)/2))
# (i) strict monotonicity, (ii) unique zero r0, (iii) asymptotics
print("Psi(0)   =", mp.nstr(Psi(0),12))        # -5.37218341923 (global min)
print("r0       =", mp.nstr(mp.findroot(Psi,6.3),16))  # 6.289835988836903
print("gamma_1  =", mp.nstr(mp.im(mp.zetazero(1)),12)) # 14.1347251417 > r0
print("Psi(g1)  =", mp.nstr(Psi(mp.im(mp.zetazero(1))),8))  # 0.81054874 > 0
for r in [20,100,1000]:                          # Psi/2pi ~ N'(r) zero density
    print(r, mp.nstr(Psi(r)/(2*mp.pi),10), mp.nstr(mp.log(mp.mpf(r)/(2*mp.pi))/(2*mp.pi),10))
\end{lstlisting}

\subsection*{A.2 Indefiniteness on primitives (Proposition~\ref{prop:indef})}
\begin{lstlisting}
import mpmath as mp
mp.mp.dps = 30
logpi = mp.log(mp.pi)
Psi = lambda r: mp.re(mp.digamma(mp.mpf('0.25')+1j*mp.mpf(r)/2)) - logpi
def Ainf_over_norm(sig):
    sig=mp.mpf(sig)
    fh=lambda r:(r**2+mp.mpf('0.25'))*mp.e**(-r**2/(2*sig**2))  # primitive: fh(+-i/2)=0
    A=(1/(2*mp.pi))*mp.quad(lambda r: fh(r)**2*Psi(r), [-mp.inf,-6.29,0,6.29,mp.inf])
    N=(1/(2*mp.pi))*mp.quad(lambda r: fh(r)**2, [-mp.inf,0,mp.inf])
    return A/N
for s in ['1','3','6','12']:
    print('sigma',s, mp.nstr(Ainf_over_norm(s),8))   # -1.82, -0.41, +0.30, +0.998
\end{lstlisting}

\subsection*{A.3 The continuum $\sltwo$ (Proposition~\ref{prop:sl2})}
\begin{lstlisting}
import numpy as np
N=6000; h=0.01; c=1.0; shift=int(round(c/h))
r=(np.arange(N)-N//2)*h
Tcont=np.diag(r)
Tc =np.zeros((N,N)); Tc[shift:,:-shift]=np.eye(N-shift)     # (T_c f)_j=f_{j-shift}
Tmc=np.zeros((N,N)); Tmc[:-shift,shift:]=np.eye(N-shift)    # (T_-c f)_j=f_{j+shift}
m=-(r**2)/c**2+r/c
H=(2.0/c)*Tcont; L=Tc; Lam=Tmc@np.diag(m)
comm=lambda A,B:A@B-B@A; s=slice(shift+20,N-shift-20); e=lambda M:np.max(np.abs(M[s,s]))
print('[Tcont,L]-cL', e(comm(Tcont,L)-c*L))    # ~3e-15
print('[H,L]-2L    ', e(comm(H,L)-2*L))         # ~7e-15
print('[H,Lam]+2Lam', e(comm(H,Lam)+2*Lam))     # ~7e-12
print('[L,Lam]-H   ', e(comm(L,Lam)-H))         # ~3e-13
\end{lstlisting}

\subsection*{A.4 Transport identity and the dichotomy comparison (Props.~\ref{prop:transport},~\ref{prop:ps})}
\begin{lstlisting}
import mpmath as mp
mp.mp.dps=25
logpi=mp.log(mp.pi)
Psi=lambda r: mp.re(mp.digamma(mp.mpf('0.25')+1j*mp.mpf(r)/2))-logpi
def vonmangoldt(nmax):
    spf=list(range(nmax+1)); i=2
    while i*i<=nmax:
        if spf[i]==i:
            for j in range(i*i,nmax+1,i):
                if spf[j]==j: spf[j]=i
        i+=1
    L=[mp.mpf(0)]*(nmax+1)
    for n in range(2,nmax+1):
        p=spf[n]; m=n
        while m%p==0: m//=p
        if m==1: L[n]=mp.log(p)
    return L
NMAX=20000; Lam=vonmangoldt(NMAX); Z=[mp.im(mp.zetazero(k)) for k in range(1,301)]
for s in ['2','3','4']:
    sig=mp.mpf(s); h=lambda r:(r**2+mp.mpf('0.25'))**2*mp.e**(-r**2/sig**2)
    g=lambda u:(1/(2*mp.pi))*mp.quad(lambda r:h(r)*mp.cos(u*r),[-mp.inf,0,mp.inf])
    Q=mp.fsum(2*h(z) for z in Z)
    A=(1/(2*mp.pi))*mp.quad(lambda r:h(r)*Psi(r),[-mp.inf,-6.29,0,6.29,mp.inf])
    gI=(1/(2*mp.pi))*mp.quad(lambda r:h(r)*abs(Psi(r)),[-mp.inf,-6.29,0,6.29,mp.inf])
    P=mp.fsum(2*(Lam[n]/mp.sqrt(n))*g(mp.log(n)) for n in range(2,NMAX+1) if Lam[n]!=0)
    print(s,'Q',mp.nstr(Q,6),'A',mp.nstr(A,6),'P',mp.nstr(P,6),
          'Q-(A-P)',mp.nstr(Q-(A-P),4),'  g_inf',mp.nstr(gI,6))
# g_inf >= 0 but differs from Q by 10^2..10^18: the grading yields g_inf, not Q.
\end{lstlisting}

\begin{thebibliography}{9}
\bibitem{RP8} D.~A.~Trejo Pizzo, \emph{A Weil--Kre\u\i n realization of the explicit formula}, preprint, 2026.
\bibitem{RP13} D.~A.~Trejo Pizzo, \emph{Two modern routes and the wrong-sign obstruction}, preprint, 2026.
\bibitem{RP15} D.~A.~Trejo Pizzo, \emph{Why the Riemann Hypothesis resists}, preprint, 2026.
\bibitem{RP16} D.~A.~Trejo Pizzo, \emph{The Pontryagin index and the finite-defect route}, preprint, 2026.
\bibitem{RP17} D.~A.~Trejo Pizzo, \emph{An unconditional finite bottom for the localized Weil form: a verified second-moment identity, an almost-everywhere clean bound, and the residue as a uniform short-interval prime variance}, preprint, 2026.

\bibitem{Weil} A.~Weil, \emph{Sur les ``formules explicites'' de la th\'eorie des nombres premiers}, Comm. S\'em. Math. Lund (1952).
\bibitem{Deninger} C.~Deninger, \emph{Some analogies between number theory and dynamical systems on foliated spaces}, ICM (1998).
\bibitem{bgWeil49} A.~Weil, \emph{Numbers of solutions of equations in finite fields}, Bull.\ Amer.\ Math.\ Soc.\ \textbf{55} (1949), 497--508.
\bibitem{bgDeligne} P.~Deligne, \emph{La conjecture de Weil. I}, Publ.\ Math.\ IH\'ES \textbf{43} (1974), 273--307.
\bibitem{bgConnes99} A.~Connes, \emph{Trace formula in noncommutative geometry and the zeros of the Riemann zeta function}, Selecta Math.\ (N.S.) \textbf{5} (1999), 29--106.
\bibitem{bgCC10} A.~Connes and C.~Consani, \emph{Schemes over $\mathbb{F}_1$ and zeta functions}, Compos.\ Math.\ \textbf{146} (2010), 1383--1415.
\bibitem{bgDeninger} C.~Deninger, \emph{Some analogies between number theory and dynamical systems on foliated spaces}, Doc.\ Math.\ ICM \textbf{I} (1998), 163--186.
\bibitem{bgBK} M.~V.~Berry and J.~P.~Keating, \emph{The Riemann zeros and eigenvalue asymptotics}, SIAM Rev.\ \textbf{41} (1999), 236--266.
\bibitem{bgMont73} H.~L.~Montgomery, \emph{The pair correlation of zeros of the zeta function}, Proc.\ Sympos.\ Pure Math.\ \textbf{24} (1973), 181--193.
\bibitem{bgKS} N.~M.~Katz and P.~Sarnak, \emph{Zeroes of zeta functions and symmetry}, Bull.\ Amer.\ Math.\ Soc.\ \textbf{36} (1999), 1--26.
\bibitem{bgBombieri00} E.~Bombieri, \emph{Problems of the Millennium: The Riemann Hypothesis}, Clay Mathematics Institute, 2000.
\bibitem{bgTitchmarsh} E.~C.~Titchmarsh, \emph{The Theory of the Riemann Zeta-Function}, 2nd ed., rev.\ D.~R.~Heath-Brown, Oxford Univ.\ Press, 1986.
\bibitem{bgEdwards} H.~M.~Edwards, \emph{Riemann's Zeta Function}, Academic Press, 1974.
\bibitem{bgIK} H.~Iwaniec and E.~Kowalski, \emph{Analytic Number Theory}, Amer.\ Math.\ Soc.\ Colloq.\ Publ.\ \textbf{53}, AMS, 2004.
\end{thebibliography}

\end{document}
